\section{Two implementation traps}\label{sec:traps}

This section describes two implementation defects. Neither prevents
convergence to the correct solution at the tolerances usually reported, so
neither is visible from final errors. Both change measured cost substantially.
We present them as evidence for the protocol of \cref{sec:protocol}. They are
not new discoveries.

\subsection{Trap 1: redundant Jacobian factorization}\label{sec:trap1}
A semismooth Newton step at $y_k$ needs the spectral decomposition of
$C(y_k)$ to apply the generalized Jacobian~\cite{qisun2006quadratically}. That
decomposition was already computed when $\theta(y_k)$ and $\nabla\theta(y_k)$
were evaluated, at the accepted trial of the previous line search. An
implementation that recomputes it pays one extra \EVD\ per outer iteration and
gains nothing. On the literature matrix \texttt{cor1399} (\cref{sec:real}),
where both variants take the same six Newton steps, recomputing the
decomposition raises Newton's cost from 7 to 13 \EVDs, an increase of about
86\%, with an identical final error
(\texttt{results/timing\_session1/timing\_real\_matched.csv}). Newton-SIN-BH caches the decomposition throughout
this paper; Newton-SIN, the naive-Armijo control of \cref{sec:trap2}, does
not, so its \EVD\ counts already include this redundant factorization on top
of the naive test's own cost. (The 56-stall finding of \cref{sec:trap2} is
about convergence, not \EVD\ counts, so this confounding does not affect it.)
The same class of defect turned up in our first
reimplementation of AGD-SDAJ (\cref{app:agd}, defect~1). It was found in the
same way, by
counting primitives one at a time against a published count.

\subsection{Trap 2: the finite-precision Armijo test}\label{sec:trap2}
An objective-value sufficient-decrease test,
\begin{equation}\label{eq:armijo}
  \theta(y_k+\alpha d_k)\le\theta(y_k)+c\,\alpha\,\nabla\theta(y_k)^{\mathsf T}d_k,
\end{equation}
becomes a plain comparison of two rounded copies of $\theta_k$ once the
predicted decrease $c\,\alpha\,|\nabla\theta(y_k)^{\mathsf T}d_k|$ falls
below half an ulp of $\theta_k$: the right-hand side then rounds to
$\theta_k$ itself, and the test degenerates to $\mathrm{fl}(\theta_t)\le
\mathrm{fl}(\theta_k)$. It stops judging the step correctly once the
achieved decrease is itself buried in that same rounding error, near
$\theta_k$. Near the solution
$|\nabla\theta^{\mathsf T}d|\gtrsim\nrm{\nabla\theta}^2$ (equality for a
gradient-type step; for Newton's step it can be considerably larger), so
this second, more severe threshold --- which concerns whether the
\emph{achieved} decrease is resolvable at all, rather than the
$c\cdot\alpha$ sufficient-decrease margin --- falls at a gradient norm of roughly
$\sqrt{u\,|\theta|}$, where $u$ is the unit roundoff. Past that point,
whether a trial passes is decided by rounding noise rather than by the
step's real merit: a step that
improves both the gradient and the forward error can still be rejected, and
the search must halve until it happens to land where the noise lets a
trial through. In this regime Newton does not halt outright; it creeps,
accepting a string of steps far smaller than warranted until it exhausts its
iteration budget without converging --- a stall in the sense used below, as
the counts just below show. None of this is new. The
cancellation problem is well known in general
optimization~\cite{hager2005new}, and Higham~\cite[Ch.~17]{higham2002accuracy}
warns more broadly that stopping tests set below what rounding allows can be
unsatisfiable even by the rounded exact solution. For the NCM dual,
Borsdorf~\cite[\S4.7.1]{borsdorf2007newton} derived the threshold (about
$10^{-8}$ in $\nrm{\nabla\theta}$ when $|\theta|\approx1$) and proposed
alternatives, which became the
modified step selection of Borsdorf and Higham~\cite[Sec.~3.3,
Alg.~4.1]{borsdorf2010preconditioned}. Their Section~3.3 observes that when
$|\phi(x)|=1$, $\nrm{p}_2<u^{1/2}$ and $\nrm{\nabla\phi(x)}_2<u^{1/2}/2$, the
computed values $fl(\phi(x+p))$ and $fl(\phi(x))$ coincide, so that the Armijo
condition cannot be verified at all. Their Algorithm~4.1, Step~6, then tests
whether the two objective values are \emph{nearly equal}, using
$|a-b|<\gamma u(1+|a|+|b|)$ with $\gamma=100$, and if so accepts the full
inexact Newton step provided
$\nrm{\nabla f(y_k+d_k)}_2/\nrm{\nabla f(y_k)}_2\le1-\mu$, falling
back to a unit steepest-descent step otherwise. Our implementation uses their
parameter values, $\mubh=0.9$, $\rho=0.5$ and $\sigma=10^{-4}$ (we write
$\mubh$ for their gradient-reduction parameter, since $\mu$ denotes the
spectrum vector of \cref{def:family}).

\paragraph{Semismooth Newton.}
With the naive test~\eqref{eq:armijo}, Newton stalls (exits at
\texttt{max\_iter} without converging) on 56 of the 270 KKT
instances; across all 270 instances it uses 151{,}471 backtracks in total.
With Borsdorf--Higham globalization
(parameters $\mubh=0.9$, $\rho=0.5$, $\sigma=10^{-4}$, and $\gamma=100$ in the
near-equality test), it converges on all 270 with 8 backtracks in total. The count of
56 was reproduced exactly in a separate timing run of the same code: at
matched tolerance, the naive-Armijo control exits at \texttt{max\_iter} on
56/270. Both counts are in the shipped result files: the 56 is the
\texttt{exit} column of the \texttt{Newton-SIN} rows of
\texttt{timing\_kkt270\_matched.csv}, and the 8 is the
\texttt{linesearch\_trials} total of the \texttt{Newton-SIN-BH} rows of that
file, which agrees with the \texttt{backtracks} column of
\texttt{ranking\_kkt270\_v2.csv.gz}. The backtrack total of 151{,}471 is the
same run's own count of rejected trials, cross-checked against its
recorded per-trial trajectory (\texttt{code/trace\_armijo\_stall.jl}). Of the 5{,}677
rejected trials of the undamped Newton step ($\alpha=1$, before any
backtracking), every one, without exception, would, if accepted, have
reduced both the gradient norm and the exact forward error; counting every
rejected trial, including those reached only after backtracking, 115{,}211
of the 151{,}471 (76\%) would have.

\paragraph{AGD-SDAJ.}
The same failure occurs in AGD-SDAJ~\cite{huynh2025accelerated}, a separately
published method with a different algorithmic structure. Line~5 of
Algorithm~1 in that paper is the objective-value test~\eqref{eq:armijo}. Run
to the common tolerance $10^{-11}$, AGD-SDAJ as published exhausts the 4000-\EVD\
budget on 18 of the 270 instances (the budget is checked at the start of each
outer iteration, so a stalled run ends at between 4000 and 4035 \EVDs). On instance
\texttt{degen-r20-m1-d1e-10-p2}, the trajectory contains 3198 rejected
monotone line-search trials, with the forward error $\fro{X_k-\Xs}$ and
$\nrm{\nabla\theta}$ fixed at $6.0935\times10^{-9}$ and
$3.5986\times10^{-9}$. The predicted decrease is about $1.3\times10^{-21}$,
far below the rounding error in $\theta$, which is about $u|\theta|\approx
7\times10^{-14}$ because $\theta(\ys)\approx3\times10^{2}$ on this instance.
We adapted the Borsdorf--Higham near-equality test to
Algorithm~1's step: when the objective difference cannot be resolved, a trial
is accepted if it reduces $\nrm{\nabla\theta}$. The test is weaker than the
Newton safeguard's $(1-\mubh)$ reduction by design: that reduction is calibrated
to Newton steps, while near the precision floor a gradient step only inches the
gradient down, so the stronger test never fires. With
this change, all 18
stalls disappear (\cref{tab:agd-stall}). Across the family, the number of
instances reaching the tolerance rises from 252/270 to 270/270, and the
worst-case cost falls from 4035 \EVDs\ to 131.

\begin{table}[t]
\centering
\caption{AGD-SDAJ on \texttt{degen-r20-m1-d1e-10-p2}, as published and with
the Borsdorf--Higham adaptation.}
\label{tab:agd-stall}
\small
\begin{tabular}{lrrlrr}
\toprule
variant & outer & \EVDs & exit & final $\nrm{\nabla\theta}$ & $\fro{X-\Xs}$\\
\midrule
AGD-SDAJ (as published) & 115 & 4001 & \EVD\ budget & $3.60\times10^{-9}$ & $6.09\times10^{-9}$\\
AGD-SDAJ-BH             & 10  & 67   & converged    & $9.40\times10^{-12}$ & $1.82\times10^{-11}$\\
\bottomrule
\end{tabular}
\end{table}

\subsection{Scope of these findings}
Three points limit what these results show.
\begin{enumerate}
  \item \emph{The phenomenon is known.} Borsdorf analysed it for the NCM
    dual~\cite[\S4.7.1]{borsdorf2007newton}, and the fix was published
    in~\cite{borsdorf2010preconditioned}. What we add is limited to (i) its
    recurrence in a separately published method that uses a naive Armijo
    test, where it is invisible under that method's native stopping rule,
    and (ii) its quantification on the controlled family (56/270 and 18/270
    stalls). For semismooth Newton the remedy is the published
    Borsdorf--Higham modified line search; for AGD-SDAJ it is our adaptation
    of their near-equality safeguard to the Algorithm~1 step of that method.
  \item \emph{It does not affect Huynh and Hwang's reported results.} Their
    stopping rule, $\nrm{\nabla\theta}_2\le10^{-7}n$, equals $10^{-5}$ on
    these instances and stops the method far above the precision floor at
    which the test fails. The stall appears only when the method is pushed to
    a tight common tolerance, which is what a matched-accuracy comparison
    requires. On their own instances, our reimplementation reproduces their
    published counts exactly for P8 and P9, and matches P7 in outer
    iterations and residual with one differing backtracking decision
    (\cref{sec:repro}, \cref{app:agd}).
  \item \emph{The adaptation is ours.} AGD-SDAJ-BH is our modification and is
    not specified by the original authors. We therefore report it alongside
    the unmodified method on every test set where both were run (all except
    the 135 higher-rank instances of \cref{sec:highrank}).
\end{enumerate}
The general point is that whether a method appears to converge at a given
tolerance can depend on a line-search detail that is unrelated to the
algorithm's design. A benchmark that stops each method under its own rule
never exposes this.
